\section{Fire and Ice: The Solar Race}

This is the first of several posts about the various races of the spirit described by \textbf{Julius Evola}: The Solar Race. Evola owes his typology of spirits of the race to \textbf{J J Bachofen}`s study of \textbf{Mother Right}, although with a creative misreading of the text, Evola extracts a deeper meaning from it.

Bachofen delineated a historical process of cultures from haeterism to matriarchy, culminating in patriarchy, which he admitted was

\begin{quotex}
the highest calling, the sublimation of earthly existence to the purity of the divine father principle.
\end{quotex}


Bachofen traces the patriarchal system in the ancient world to Athens and to the Roman Imperium. In them, there was Solarity.

As we have seen, however, Evola rejects any notion of ``progress'' in history. Rather he sees a degeneration. Certain ancient peoples were feminine in nature, and hence endured various matriarchal systems as described by Bachofen. On the other hand, the Hyperborean, or Nordic, peoples, were originally masculine, or equivalently Solar. Subsequent racial degeneration has led to the situation of the present era.

Although Evola gives us a very brief outline, he entire oeuvre is dedicated to the description of the Solar race. Fire and Ice are the paradoxically conjoined symbols of the Solar qualities. The Solar beings are described by Evola as being in this world, but not of it: they feel themselves as strangers in a strange land. Only ignorance or a waking ``sleep'' keep them bound to specific modes of being in the world. Gnosis, or awakening, what Guenon calls the Supreme Identity, liberates them from attachment to the world. Evola claims that such doctrines of awakening can only apply to the Solar race; attempts to extend such ideas as universal can lead only to misunderstandings.

This series, we hope, should make clear that the problem of the modern world is not a problem of thinking but rather of being. The error of the intellectual type is the belief that thinking the correct thoughts will somehow rectify things. Hence, there is incessant chatter and debate, which, apart from a few random conversions, always ends up in aporia.

The fundamental thing that matters is a change in the level of being.

\paragraph{Text}


The doctrine of the third degree of race must fundamentally limit its researches to the sphere of influence of a given race of the spirit and of its primordial tradition, following its developments, mutations (paravariations), and also distortions in the cycle that coincide with them and in which they act and react in the confrontation with the influences of different races or new environmental conditions. Once research is circumscribed that way, we arrive at a more limited concept of race, corresponding to that of the various differentiations or articulations of the primary element of such a cycle. Along the same lines, we cannot of course conceive of an atomic separation of the various ``races of the spirit'': their differences do not exclude relationships of descent or even of different hierarchical rank.

We previously sketched out the typology of races of the spirit, as far as it concerns the specific human cycle of the hyperborean race, in the second part of our work \emph{Revolt against the Modern World} (with special regard to the properly traditional and spiritual aspect), and in the selection of the writings of J. J. Bachofen, and in a re-interpretation of them in a racial sense, included in the volume we edited under the title \emph{La razza solare}. For a fuller treatment, the reader is therefore referred back to these two works. Here we will give only a brief schematic synthesis, necessarily lacking any justifying elements.

The solar or Olympic-solar race, corresponding to the hyperborean lineage and tradition, is to be considered superior and anterior to all the others in the cycle in question. It has as characteristic a type of ``natural supernaturalism''; spirit and power, calm dominance, and readiness for precise and absolute action, a sense of ``centrality'' and of ``firmness'', and in its exterior effects, that virtue to which the ancients attributed a ``numinous'' quality (from \emph{numen}), as superiority that imposes itself directly and irresistibly, that awakens simultaneously both fear and admiration, are marks of this ``race of the spirit'', by means of which it is naturally predestined to command and, at the limit, to the regal function. Fire and ice are united in it as in the ambiguous symbols of the original Nordic homeland and of the cycle in which it is eminently and primordially manifested: ice, as transcendence and inaccessibility; fire, as the specifically radiant solar quality of the men who give life, who awaken life, and generate light, but always from a supreme distance and almost with indifference, like the wake of a boat, not through enthusiasm, inclination, or human concern.

The ancient symbol of gold has always had connections with this form of spirituality. In the political forms of the origins, it acted as the substrate to sacred, or divine, regality, i.e., to the union of the two powers, of the regal and priestly function, the latter understood in a higher sense that will soon be made clear. The symbolic designations of the ``divine'' or ``celestial'' race must refer to the absence of the dualistic sentiment in relations to supernatural reality, something that however is quite distinct from everything that is immanence or Promethean arrogance in the modern sense: it is not about men who believe they are gods, but of the type of men who naturally, through a not yet dim memory of origins and a condition of soul and body that do not cripple such a memory, feel they do not properly belong to the terrestrial race, to the extent of believing themselves men only by accident, or through ``ignorance'', or ``sleep''. The two terms
\emph{vidya} and \emph{avidya} of the ancient Indo-Aryan teaching, that mean respectively ``knowledge'' (of the ``supreme identity'') and ``ignorance'' (that leads to the self-identification with one of the forms or modes of being of the conditioned world), are to be understood exactly in this context. When these terms are attributed to a different human condition and a different race of the spirit, that is, when turned into ``philosophical'' terms, they lose all meaning and give rise to misunderstandings of various types.


\flrightit{Posted on 2015-05-31 by Cologero}

\begin{center}* * *\end{center}

\begin{footnotesize}
\begin{sffamily}
\texttt{Mark Citadel on 2015-06-03 at 11:20 said:}

Interesting how today's Nordics are some of the most spiritually degenerated in the world. Does the fish rot from the head, or is there a special historical sociohistorical reason for the total implosion in Tradition for countries like Sweden and Norway?

\hfill

\texttt{obscure on 2015-06-03 at 18:18 said:}

Mark,

One the one hand, there is the Protestant tendency against any established authority or hierarchy; a tendency which prevents a community from being more than incidental. On the other hand, there is the current incidental cause operating on Northern Europeans in general i.e. post-WW2 race-guilt. These are the two general causes which are readily conceivable. One may also note how easily a certain self-concept may be diffused throughout a sparse population.

Lastly, wherever a hierarchy cannot be established which actualizes otherwise abstract principles, then there will be a tendency towards some form of political hedonism. Such hedonism may be quasi-moral and egalitarian or individualistic and egoistic, but it certainly won't be representative of true principles except accidentally i.e. insofar as all reality is dependent upon some degree of truth. This last observation should mean that there are in fact `seeds' present in the North even though the ground appears barren.

\hfill

\texttt{aegishjalmer on 2015-06-19 at 23:14 said:}

Also, let's not forget that the Nordics were the last to be Christianised generally in Europe (c. 1000 in Iceland) and this was of limited success and depth. Hence they were quick to embrace an outlook that was anti Roman (Protestantism) and eventually, modern paganism in various ways, such as Asatru and general modern hedonism. The Nordics also have a strong culture of individualism in many ways, and are very transparent and open to other cultures historically, both in their exploration and contact, and in allowing others in. I don't think these ideas cover all there is to it, but national culture is one aspect.

\hfill

\texttt{Aleksandar Gueric on 2015-10-16 at 11:18 said:}

I just received a book ''Among the Hyperboreans'' by Miloš Crnjanski. This interesting autobiography consists of a series of ruminations and meditations clothed in multiform melancholy on Europe at the end of 1930s and early 1940s, when he worked as a diplomat in Rome.

When Miloš Crnjanski, one of the greatest Serbian poets and authors of the last century, visited Nordic countries hoping to see a miracle of mid 20th century Socialism -- in Denmark, Sweden and Norway, he realized that there was all along something else which brought him there. It was the North itself. So he headed further, to the shores of the North sea, wanting to dissolve in its waters. In his own words, '\,'something'\,' changed all those who searched for the North, and he wants to discover it's secrets.

His life embodied that '\,'stranger in a strange land'\,'concept. Exiled from communist Yugoslavia after the war, he spent life journeying through Europe, examining events, people and his own beliefs -- finding errors in general ideas of civil rights, liberalism and democracy, who only began to take forms in that time, while carefully waiting his return to Ithaca. Little known fact is that he supported Musolini.

After the fall into despair and melancholy caused by the sickness of every day reality, what is left for us than to search for that Ideal, where Love, Death, Life, Dream and Hope unite? Will you struggle in the mud, and, if you are aninmal in human clothing, you certainly will enjoy it, or will you reach for the stars? For the Absolute, Sacred Center?

Milos' travels gravitated around two poles -- Rome and Scandinavia, horizontally; vertically -- between Hyperborea and Sumatra, which I, mythologicaly, recognize as Lemuria. This is Solar and Lunar dichotomy expressed often in poet's tendencies towards abstract Lunar spheres of the mind. He was, afterall, still a Romanticist, thus all the Moon, autumn and death imagery in his writings. But behind poetic expressionism, his was of a Solar race. Many Solar individuals mantle different personas. Romanticism and it's ideals, for example, often lie at the beginning of a path in those born within a Solar race. It's an early, raw stage of awakened Primordial Man. Ecstasis of war, death and Platonic love is truly titanic, but not Olympic, so it often entraps those desparate souls at the very beginning of their search for a way out of the Mundane into the Golden. This is why so many Romantics celebrated war and even took part in some of them, from Spain to Greece, seeking heroic death in the Age when heroic death (or life) is almost unknown. This is the echo of their Solar origin. Even Novalis, archetypal Romantic, who venerated Sofia and Mother Night, once said that '\,'Man is a Sun'\,'. Little did he know (or perhaps he did), that this Man is quite rare. Solar Man is a man of Tradition, one and Perennial, and Novalis, as a man of Tradition, saw in French Revolution death of Europe.

Friedrich Hölderlin was a priest of Helios. This is a common Romantic naivette. If only he went further along the path of Solar Man, he would ascend to a priesthod of Apollo.

Crnjanski, like Evola, felt at home both in Rome and in Nordic countries. Eternal Mother Amor, the ever nourishing Fountain and distant North, Solar and Virile heroic coldness, parents of the European, Aryan Soul. Evola, on the other hand, felt no affection (unlike Guenon) with the Sumatrian -- or Lunar center.

A symptom each Solar Man, a '\,'stranger in a strange land'\,' will share is subtle Melancholy, expressed through a desire, deep rooted yearning towards the Absolute, while world is seen through the eyes of mono no aware. This isn't some pathetic sadness and infantile escapism but (supra)natural effect of a Soul awareness elevated from the Animal and connected to the Higher Center. It pulls towards the Transcendental, above the mass and material. Gravity of this Higher Center can break the Promethean chains, but only through the endeavor worthy of Herakle himself, meaning we must already have a divine potentiality. What happens when Olympian Fire is shared with everyone, both worthy and unworthy, we know through the myth of Prometheus. Sadly, modern Europeans, even some calling themsleves Traditionalist or '\,'of the Right'\,', follow the Promethean path.

Heraklean labors, just like Odysseus' journey, are part of the Heroic, Solar Path. I don't share common opinion that we all inherently have our own, personal Odyssey, or, carry our own cross or hang from Yggdrassil. Or, that we are all immortal souls. That comes later, after accepting a call and undergoing a path. Most people don't follow any path but just pave a road.

In the end, all that's left is to wander these mysterios roads and gravitate towards different spiritual centers, according to our curent inner balance of psychic energies.

\hfill

\texttt{Tyler on 2018-05-22 at 12:59 said:}

Why is the moon masculine and the sun feminine in Germanic mythology?

\hfill

\texttt{Dark Worship on 2020-01-14 at 18:51 said:}

I believe the ''high culture of individualism'' is when people seek easy going and mediocre life(i for exemple, think in this way). It's not every person that has enough will to go to a war, a guerrila warfare etc People will die on the same way, being killed by migrants on street or being killed in war. There's no escape, i dont believe in superhumans. Maybe the individualists just want peace being locked in home.

\hfill

\texttt{Alice on 2021-11-05 at 08:10 said:}

While Evola tries to differentiate the concept of the Promethean and the Solar Race, the latter comes across as the Anti-Christ and the epitome of what Dostoevsky calls Man-Godism. Human beings can never approach God's level because God is the Supreme Person and we exist as His servants, whereas Evola derides Bhakti as inferior and ``lunar.'' This doesn't seem any different from Theosophy or other New Age religions.

\hfill

\texttt{S.P. on 2021-11-05 at 10:54 said:}

I disagree, although it is a common confusion. What Evola mentions does not consist in the identification of the human being as an individual man with the supreme God, but rather the self-knowledge / realization of the I, which implies the overcoming of the human state and therefore of the egoic individuality.

\hfill

\end{sffamily}
\end{footnotesize}

